\section{Conclusion}
In this paper, we proposed Diff-LM-GZSL, a novel language-driven latent diffusion framework for continuous semantic alignment in zero-shot fault diagnosis. The proposed framework effectively addresses the limitations of traditional attribute-based zero-shot learning by leveraging the rich contextual knowledge of large language models combined with the superior generative capabilities of latent diffusion probabilistic models. 

The main contributions of this work are three-fold: 
First, we introduced LLM-generated linguistic embeddings to provide richer and more discriminative semantic information compared to traditional binary attributes. By utilizing BERT-based encoders to process expert-level textual descriptions, the model captures nuanced physical characteristics that are often lost in rigid attribute matrices. 
Second, we developed a conditional latent diffusion module that synthesizes high-quality fault features by learning the conditional distribution of signal spectrograms. This module demonstrates a remarkable ability to model the high-dimensional manifold of industrial sensor data, significantly reducing the class-bias inherent in generalized zero-shot learning. 
Third, we implemented a semantic-consistent alignment loss based on Maximum Mean Discrepancy (MMD) to ensure the generated features are physically meaningful and domain-invariant, bridging the gap between semantic descriptions and raw sensor measurements.

Extensive experiments conducted on the TEP, Hydraulic System, and CWRU bearing datasets demonstrate that Diff-LM-GZSL consistently achieves state-of-the-art performance. The experimental findings revealed significant improvements in the Harmonic Mean ($H$) across all three benchmarks, with gains of 4.2\%, 4.9\%, and 6.8\% respectively over the strongest generative baselines. These results underscore the core advantage of language-driven semantic spaces over binary attributes, particularly in complex industrial scenarios where fault symptoms are subtle and multi-modal.

Future research will focus on several key directions. We plan to explore multi-modal fusion by incorporating acoustic and thermal imaging data alongside vibration signals to further enhance diagnostic reliability. Additionally, optimizing the diffusion sampling process via consistency models or distillation techniques is a priority to enable real-time inference and faster adaptation in industrial edge computing environments. Finally, we aim to extend this framework to few-shot and incremental learning settings to better handle the evolving nature of industrial fault patterns.
